% ============================================================
%  Глава 5 — Линеаризация и дискретизация
% ============================================================
\chapter{Из нелинейной ОДУ в матрицы $\mat{A}_d, \mat{B}_d$}
\label{ch:linearization}

\section{Вопрос}

Нелинейная модель \eqref{eq:nonlinear_ode} --- честная, но \textbf{нерешаемая
аналитически} и плохо подходящая для оптимизаторов. ЛКР и MPC, которые
мы будем строить дальше, работают \emph{только с линейными моделями}
в дискретном времени:

\[
  \vect{x}_{k+1} = \mat{A}_d\,\vect{x}_k + \mat{B}_d\,\vect{u}_k.
\]

Надо превратить \eqref{eq:nonlinear_ode} в эту форму.

\begin{ideabox}
Две операции, обе из учебника Козлова:
\begin{enumerate}
  \item \textbf{Линеаризация.} Разложение в ряд Тейлора первого порядка
        вокруг \emph{опорной точки} (равномерное движение вперёд).
        Получаем \emph{непрерывные} матрицы $\mat{A}, \mat{B}$.
  \item \textbf{Дискретизация.} Переход от $\dot{\vect{x}} = \mat{A}\vect{x}
        + \mat{B}\vect{u}$ к $\vect{x}_{k+1} = \mat{A}_d\vect{x}_k +
        \mat{B}_d\vect{u}_k$ через ZOH (zero-order hold).
\end{enumerate}
В этой главе пройдём по обеим.
\end{ideabox}

\section{Выбор опорной точки}

\textbf{Опорная точка} (operating point, англ.\,\emph{trim point}) --- это
тот режим, вокруг которого мы линеаризуем. Выберем его так, чтобы он
был \emph{реальной целью} сценария: робот едет вперёд с постоянной
скоростью $v_0$, ничего не вращая.

\begin{equation}
  \bar{\vect{x}} = (\bar{p}_x,\; \bar{p}_y,\; \underbrace{0}_{\bar\theta},\;
                    v_0,\; \underbrace{0}_{\bar\omega})\T,
  \qquad
  \bar{\vect{u}} = (v_0,\; 0)\T.
  \label{eq:op_point}
\end{equation}

Подставим $\bar{\vect{x}}, \bar{\vect{u}}$ в \eqref{eq:nonlinear_ode}:

\[
  \vect{f}(\bar{\vect{x}}, \bar{\vect{u}}) =
  (v_0\cos 0,\, v_0\sin 0,\, 0,\, (v_0 - v_0)/\tau_v,\, (0-0)/\tau_\omega)\T
  = (v_0, 0, 0, 0, 0)\T.
\]

Координата $p_x$ растёт со скоростью $v_0$ --- логично; всё остальное
не меняется. Опорная точка корректна: при $\vect{x} = \bar{\vect{x}}$,
$\vect{u} = \bar{\vect{u}}$ робот действительно равномерно едет
вперёд.

\section{Ряд Тейлора}

Пусть $\Delta\vect{x} = \vect{x} - \bar{\vect{x}}$,
$\Delta\vect{u} = \vect{u} - \bar{\vect{u}}$ --- \emph{малые}
отклонения от опорной. Разложим $\vect{f}$ в ряд:

\[
  \vect{f}(\vect{x}, \vect{u}) \;=\;
    \underbrace{\vect{f}(\bar{\vect{x}}, \bar{\vect{u}})}_{\dot{\bar{\vect{x}}}}
    +
    \underbrace{\left.\frac{\partial \vect{f}}{\partial \vect{x}}\right|_*}_{\mat{A}}
    \Delta\vect{x}
    +
    \underbrace{\left.\frac{\partial \vect{f}}{\partial \vect{u}}\right|_*}_{\mat{B}}
    \Delta\vect{u}
    + O(\|\Delta\|^2).
\]

Вычитая $\dot{\bar{\vect{x}}}$ из обеих частей и обозначив
$\Delta\dot{\vect{x}} = \dot{\vect{x}} - \dot{\bar{\vect{x}}}$, отбрасывая
квадратичные члены, получаем:

\begin{formulabox}[title={Линеаризованная модель}]
\begin{equation}
  \Delta\dot{\vect{x}}(t) \;=\; \mat{A}\,\Delta\vect{x}(t)
                                 + \mat{B}\,\Delta\vect{u}(t).
  \label{eq:linear_ode}
\end{equation}
\end{formulabox}

\textbf{Важно:} формула \eqref{eq:linear_ode} верна \emph{только в окрестности}
$\bar{\vect{x}}, \bar{\vect{u}}$. Если робот сильно отклонится от опорной
траектории (например, разверулся на $90^\circ$), линейная модель станет
заметно врать.

\subsection{Вычисление матрицы $\mat{A}$}

Матрица Якоби $\mat{A} = \partial \vect{f}/\partial \vect{x}$ --- это
матрица $5 \times 5$ со столбцами $\partial \vect{f}/\partial x_j$.
Считаем частные производные функций $f_1, \ldots, f_5$ по $x_j \in \{p_x,
p_y, \theta, v, \omega\}$ и подставляем опорные значения
$\bar{\theta} = 0$, $\bar{v} = v_0$.

\textbf{Строка 1:} $f_1 = v\cos\theta$.
\[
  \frac{\partial f_1}{\partial p_x} = 0,
  \quad
  \frac{\partial f_1}{\partial p_y} = 0,
  \quad
  \frac{\partial f_1}{\partial \theta}\bigg|_* = -v\sin\theta\big|_0 = 0,
  \quad
  \frac{\partial f_1}{\partial v}\bigg|_* = \cos\theta\big|_0 = 1,
  \quad
  \frac{\partial f_1}{\partial \omega} = 0.
\]

\textbf{Строка 2:} $f_2 = v\sin\theta$.
\[
  \frac{\partial f_2}{\partial \theta}\bigg|_* = v\cos\theta\big|_0 = v_0,
  \quad
  \frac{\partial f_2}{\partial v}\bigg|_* = \sin\theta\big|_0 = 0.
\]
Остальные --- нули.

\textbf{Строка 3:} $f_3 = \omega$.
$\partial f_3/\partial \omega = 1$. Остальные --- нули.

\textbf{Строка 4:} $f_4 = (u_v - v)/\tau_v$.
$\partial f_4/\partial v = -1/\tau_v$. Остальные --- нули.

\textbf{Строка 5:} $f_5 = (u_\omega - \omega)/\tau_\omega$.
$\partial f_5/\partial \omega = -1/\tau_\omega$. Остальные --- нули.

Собираем:
\begin{formulabox}[title={Матрица $\mat{A}$ (непрерывная)}]
\begin{equation}
  \mat{A} = \begin{pmatrix}
    0 & 0 & 0   & 1         & 0          \\
    0 & 0 & v_0 & 0         & 0          \\
    0 & 0 & 0   & 0         & 1          \\
    0 & 0 & 0   & -1/\tau_v & 0          \\
    0 & 0 & 0   & 0         & -1/\tau_\omega
  \end{pmatrix}.
  \label{eq:A_continuous}
\end{equation}
\end{formulabox}

\subsection{Вычисление матрицы $\mat{B}$}

Аналогично для $\mat{B} = \partial \vect{f}/\partial \vect{u}$, размер $5 \times 2$:

\[
  \frac{\partial f_4}{\partial u_v} = \frac{1}{\tau_v},
  \quad
  \frac{\partial f_5}{\partial u_\omega} = \frac{1}{\tau_\omega}.
\]
Остальные элементы --- нули.

\begin{formulabox}[title={Матрица $\mat{B}$ (непрерывная)}]
\begin{equation}
  \mat{B} = \begin{pmatrix}
    0          & 0              \\
    0          & 0              \\
    0          & 0              \\
    1/\tau_v   & 0              \\
    0          & 1/\tau_\omega
  \end{pmatrix}.
  \label{eq:B_continuous}
\end{equation}
\end{formulabox}

Если сравнить эти матрицы с кодом \filepath{state\_space\_model.py:29}
(глава~\ref{ch:dynamics}, листинг) --- буквально одно и то же.
\emph{Это и есть наша модель.}

\subsection{Численный пример}

При $v_0 = 0.2$~м/с, $\tau_v = 0.15$~с, $\tau_\omega = 0.10$~с:

\begin{equation}
  \mat{A} = \begin{pmatrix}
    0 & 0 & 0    & 1      & 0 \\
    0 & 0 & 0.2  & 0      & 0 \\
    0 & 0 & 0    & 0      & 1 \\
    0 & 0 & 0    & -6.667 & 0 \\
    0 & 0 & 0    & 0      & -10
  \end{pmatrix},
  \quad
  \mat{B} = \begin{pmatrix}
    0      & 0 \\
    0      & 0 \\
    0      & 0 \\
    6.667  & 0 \\
    0      & 10
  \end{pmatrix}.
  \label{eq:AB_numeric}
\end{equation}

\section{Управляемость}

Перед тем как строить регулятор, надо убедиться, что робот \emph{в принципе}
может быть приведён из любого состояния в любое другое. Это свойство
\textbf{управляемости}.

\begin{formulabox}[title={Критерий Калмана}]
Объект полностью управляем тогда и только тогда, когда матрица
управляемости имеет полный ранг:
\begin{equation}
  \mat{S}_y = \big[\,\mat{B}\,\big|\,\mat{A}\mat{B}\,\big|\,\mat{A}^2\mat{B}
              \,\big|\,\ldots\,\big|\,\mat{A}^{n-1}\mat{B}\,\big],
  \qquad
  \rank \mat{S}_y = n.
  \label{eq:controllability}
\end{equation}
\end{formulabox}

Для нашей $n = 5$, $r = 2$: $\mat{S}_y$ имеет размер $5 \times 10$.
Численная проверка в MATLAB (\filepath{matlab/analyze\_system.m}) даёт
$\rank \mat{S}_y = 5$ при $v_0 \ne 0$ --- система полностью управляема.

\begin{warnbox}
\textbf{При $v_0 = 0$} опорная точка вырождена. В \eqref{eq:A_continuous}
второй столбец, отвечающий за связь $\dot{p}_y$ с $\theta$, обнуляется ---
координата $p_y$ становится неуправляемой через малые отклонения
$\Delta\theta$. Физически: стоя на месте, робот не умеет «поехать вбок»,
не повернувшись сильно (а это уже не «малые отклонения», и линейная
модель не работает).

\textbf{Следствие:} линеаризуем только вокруг $v_0 > 0$. Чтобы стартовать
с нуля, мы используем \emph{два разных режима}: TURN на стоящем роботе
управляется тем же MPC, но \emph{с другой опорной} (об этом --- глава
\ref{ch:turn_phase}).
\end{warnbox}

\section{Дискретизация: ZOH}

MPC работает с дискретной моделью, потому что:

\begin{itemize}
  \item Цикл управления --- дискретный (тик $T_s$).
  \item Команды моторам тоже подаются дискретно: «на интервале $[k T_s,
        (k+1) T_s)$ команда $\vect{u}_k$, потом следующая».
\end{itemize}

\begin{ideabox}
\textbf{ZOH (zero-order hold).} «Нулевого порядка удержание»: команда
$\vect{u}_k$ \emph{постоянна} на всём интервале $[k T_s, (k+1)T_s)$.
Это точно соответствует тому, как мы кормим $\texttt{cmd\_vel}$
моторам --- между тиками значение не меняется.

В этом предположении уравнение $\dot{\vect{x}} = \mat{A}\vect{x} +
\mat{B}\vect{u}$ можно проинтегрировать \emph{точно} от $kT_s$ до
$(k+1)T_s$ --- получим формулу с матричной экспонентой.
\end{ideabox}

\subsection{Аналитический вывод}

Решение линейной ОДУ с постоянной $\vect{u}_k$ на $[t_0, t_0 + T_s]$:

\[
  \vect{x}(t_0 + T_s)
    = e^{\mat{A} T_s}\,\vect{x}(t_0)
    + \int_0^{T_s} e^{\mat{A}(T_s - \tau)}\mat{B}\,d\tau \cdot \vect{u}_k.
\]

Сменой переменной $\sigma = T_s - \tau$ интеграл превращается в

\[
  \int_0^{T_s} e^{\mat{A}\sigma}\,d\sigma \cdot \mat{B}.
\]

Поэтому:

\begin{formulabox}[title={ZOH-дискретизация}]
\begin{equation}
  \boxed{\;
    \mat{A}_d = e^{\mat{A} T_s},
    \qquad
    \mat{B}_d = \left(\int_0^{T_s} e^{\mat{A}\sigma}\,d\sigma\right) \mat{B}.
    \;}
  \label{eq:zoh}
\end{equation}
\end{formulabox}

\subsection{Блочная экспонента (как считают в коде)}

Численно интеграл считать неудобно --- есть \textbf{трюк с блочной
матричной экспонентой}. Соберём блочную матрицу

\[
  \mat{M} =
  \begin{pmatrix}
    \mat{A} & \mat{B} \\
    \mat{0} & \mat{0}
  \end{pmatrix} \in \Real^{(n+r) \times (n+r)}.
\]

Аккуратное матричное возведение в экспоненту даёт:

\begin{equation}
  \exp(\mat{M} T_s) =
  \begin{pmatrix}
    e^{\mat{A} T_s} & \int_0^{T_s} e^{\mat{A}\sigma}d\sigma \cdot \mat{B} \\
    \mat{0}         & \mat{I}_r
  \end{pmatrix}
  =
  \begin{pmatrix}
    \mat{A}_d & \mat{B}_d \\
    \mat{0}   & \mat{I}_r
  \end{pmatrix}.
  \label{eq:block_exp}
\end{equation}

То есть: \textbf{одна} операция \texttt{expm} даёт сразу обе матрицы.

\begin{codebox}[title={\filepath{pi\_nodes/control/state\_space\_model.py:48}}]
\begin{lstlisting}[language=Python]
def zoh_discretize(A, B, Ts):
    """Exact ZOH discretization via block matrix exponential."""
    n, r = A.shape[0], B.shape[1]
    M = np.zeros((n + r, n + r))
    M[:n, :n] = A
    M[:n, n:] = B
    M_d = expm(M * Ts)
    Ad = M_d[:n, :n]
    Bd = M_d[:n, n:]
    return Ad, Bd
\end{lstlisting}
\end{codebox}

\subsection{Полученная дискретная модель}

\begin{formulabox}[title={Дискретная модель робота}]
\begin{equation}
  \boxed{\;\vect{x}_{k+1} = \mat{A}_d\,\vect{x}_k + \mat{B}_d\,\vect{u}_k.\;}
  \label{eq:discrete_model}
\end{equation}
\end{formulabox}

\textbf{Это финальная форма}, с которой будут работать ЛКР, MPC, симулятор.

\section{Собственные значения}

Собственные значения говорят, как себя ведёт открытая (без управления)
система.

\begin{notebox}
Связь между непрерывным $\lambda$ и дискретным $\tilde{\lambda}$:
\[
  \tilde{\lambda}_i = e^{\lambda_i T_s}.
\]
Устойчивость:
\begin{itemize}
  \item непрерывная: $\operatorname{Re}\lambda_i < 0$ (левая полуплоскость);
  \item дискретная: $|\tilde{\lambda}_i| < 1$ (внутри единичного круга).
\end{itemize}
\end{notebox}

Для нашей $\mat{A}$ \eqref{eq:A_continuous} собственные значения легко
считать «руками» --- $\mat{A}$ блочно-треугольная:

\[
  \det(\mat{A} - \lambda \mat{I}) =
  (-\lambda)^3 \cdot (-1/\tau_v - \lambda) \cdot (-1/\tau_\omega - \lambda).
\]

\begin{equation}
  \lambda_{1,2,3} = 0,
  \qquad
  \lambda_4 = -1/\tau_v,
  \qquad
  \lambda_5 = -1/\tau_\omega.
  \label{eq:eigenvalues}
\end{equation}

Три нуля --- это \emph{интеграторы}: $p_x, p_y, \theta$ накапливаются
во времени без внутреннего «возврата к нулю». Два отрицательных корня
--- инерционность моторов: возмущение по скорости затухает.

В дискретной форме: $\tilde\lambda_{1,2,3} = 1$ (на границе устойчивости),
$\tilde\lambda_4 = e^{-T_s/\tau_v} \approx 0.717$,
$\tilde\lambda_5 = e^{-T_s/\tau_\omega} \approx 0.607$.

\begin{warnbox}
\textbf{Объект не асимптотически устойчив.} Три собственных значения
на единичной окружности означают: без управления отклонения по координатам
и курсу не затухают (что логично --- робот не «возвращается домой» сам
по себе). Это норма для движущихся объектов. ЛКР/MPC сделают замкнутую
систему устойчивой.
\end{warnbox}

\section{Каноническая МПС-форма $(s, v, \theta, \omega, e_\text{int})$}

Выше мы работали с легаси-формой $(p_x, p_y, \theta, v, \omega)$,
потому что её удобно линеаризовать «с нуля» из физики. В коде
\filepath{mps\_node.py} используется \emph{другая}, каноническая форма
$(s, v, \theta, \omega, e_\text{int})$. Концептуально это та же модель,
просто:

\begin{itemize}
  \item вместо двух декартовых $(p_x, p_y)$ --- одна продольная $s$
        (мы заранее знаем, что в сценарии $p_y \equiv 0$);
  \item плюс интегральный член $e_\text{int}$, чтобы устранить
        установившуюся ошибку скорости.
\end{itemize}

Конкретные матрицы $\mat{A}, \mat{B}$ для этой формы синтезируются
в \filepath{matlab/main.m} и экспортируются в \filepath{config.yaml},
секция \texttt{mps.matrices.\{A,B,C,D\}}. Размер тот же ---  $5 \times 5$ и $5
\times 2$, но числа другие. Период тоже \emph{другой}: $T_s = 0.02$~с
(50~Гц), а не $0.05$~с легаси-контура.

\begin{notebox}
В коде эти матрицы приходят в \emph{непрерывной} форме (контракт
\filepath{docs/mps/api.md}). \texttt{mps\_node} ZOH-дискретизирует их
при загрузке/Apply --- так что одна и та же $\mat{A}$ работает
и в симе на ноуте, и на Pi, в любых обновлениях.
\end{notebox}

\section{Что мы получили}

\begin{itemize}
  \item \textbf{Линеаризация}: ряд Тейлора 1-го порядка вокруг опорного
        $(v_0, 0)\T$. Получились $\mat{A} \in \Real^{5\times 5}$,
        $\mat{B} \in \Real^{5\times 2}$ с явными формулами \eqref{eq:A_continuous},
        \eqref{eq:B_continuous}.
  \item Объект \textbf{полностью управляем} при $v_0 \ne 0$.
  \item \textbf{Дискретизация ZOH}: $\mat{A}_d = e^{\mat{A}T_s}$, $\mat{B}_d$
        --- через интеграл. Вычисляется одним вызовом \texttt{expm} от
        блочной матрицы.
  \item Собственные значения непрерывной системы: три нуля (интеграторы
        по $p_x, p_y, \theta$) и два отрицательных. Открытый объект не
        асимптотически устойчив, но управляем.
  \item Каноническая МПС-форма --- та же модель с $(s, \cdot, \theta, \omega,
        e_\text{int})$ вместо $(p_x, p_y, \theta, v, \omega)$ и
        $T_s = 0.02$~с.
\end{itemize}

\medskip
\textbf{Теперь у нас есть всё, что нужно для регулятора.} Дальше начинаем
строить управляющий контур.
